\documentclass[aspectratio=169,11pt]{beamer}

\usetheme{Madrid}
\usecolortheme{default}
\usefonttheme{professionalfonts}
\setbeamertemplate{navigation symbols}{}
\setbeamertemplate{footline}[frame number]

\usepackage{amsmath, amssymb, booktabs}

\title{Contracts, Incentives, Screening, And Moral Hazard}
\subtitle{Contract Design Inside Employment Relationships}
\author{Special Topic 7: Labor Market Design, Contracting, and Mechanism Design -- Week 3}
\date{}

\begin{document}

\begin{frame}
  \titlepage
\end{frame}

\begin{frame}{Central question}
\begin{itemize}
    \item How should labor contracts be designed when effort, type, commitment, and multitasking are imperfectly observed?
    \item Week 3 moves inside the employment relationship after the worker-firm match is formed.
    \item The labor objects are pay, performance measures, probation, promotion, monitoring, subjective ratings, retention, and risk.
    \item The research standard is translation: hidden object $\rightarrow$ observed margin $\rightarrow$ institutional variation $\rightarrow$ welfare.
\end{itemize}
\end{frame}

\begin{frame}{Principal-agent baseline}
\small
\[
  y = f(e,\theta,x) + \varepsilon
\]
\begin{itemize}
    \item Worker effort $e$ is costly and imperfectly observed.
    \item Output $y$ depends on effort, hidden type $\theta$, job design $x$, and noise $\varepsilon$.
    \item A contract conditions pay on an imperfect signal: output, attendance, quality, rating, or promotion score.
    \item Stronger incentives raise returns to effort but shift risk to workers.
    \item The core constraints are participation and incentive compatibility.
\end{itemize}
\[
  \max_{a,b} \; \mathbb{E}[y-w(y)]
  \quad \text{s.t. participation and incentive compatibility.}
\]
\end{frame}

\begin{frame}{Screening and hidden type}
\begin{itemize}
    \item Hidden type is not hidden effort.
    \item Type can mean productivity, reliability, patience, promotability, collaboration, risk, or self-control.
    \item Screening devices: probation, promotion ladders, nonlinear bonuses, referral rules, benefit menus, and schedule flexibility.
    \item A contract change can alter who enters, who exits, and how continuing workers behave.
    \item Lazear's performance-pay setting is the teaching benchmark because productivity gains combine incentives and sorting.
\end{itemize}
\end{frame}

\begin{frame}{Incentives and measured productivity}
\small
\begin{tabular}{p{0.25\linewidth} p{0.33\linewidth} p{0.31\linewidth}}
\toprule
Contract object & Observable labor margin & Interpretation risk \\
\midrule
Piece rate & output per hour, pay, quits & effort versus sorting \\
Team bonus & store output, peer output & free riding versus peer pressure \\
Manager incentive & worker allocation, output, inequality & effort versus favoritism \\
Goal or prize & task completion, quality, persistence & motivation versus gaming \\
\bottomrule
\end{tabular}
\vspace{0.25cm}
\begin{itemize}
    \item Output is useful, but it is not effort itself.
    \item Good papers ask whether output rose through effort, time, sorting, task reallocation, learning, or gaming.
\end{itemize}
\end{frame}

\begin{frame}{Multitasking, subjectivity, and relational incentives}
\small
\[
  y_j = f_j(e_j) + \varepsilon_j,
  \qquad j=1,\dots,J.
\]
\begin{itemize}
    \item If only some tasks are measured, workers reallocate effort toward rewarded tasks.
    \item Subjective evaluation can recover soft information: cooperation, judgment, mentoring, quality, and reliability.
    \item Subjectivity also creates influence activities and evaluator-specific effort.
    \item Relational incentives use repeated interaction, discretionary bonuses, promotion promises, and reputation.
    \item The labor question is which signal is rewarded and which valuable work becomes invisible.
\end{itemize}
\end{frame}

\begin{frame}{Commitment and worker-side demand}
\begin{itemize}
    \item Some contracts solve worker-side commitment problems, not only firm-side agency problems.
    \item Workers may value targets, schedules, deadlines, or penalties that help them follow through on desired effort.
    \item Contract take-up can reveal commitment demand, but it can also reflect risk preferences, liquidity constraints, or outside options.
    \item Welfare interpretation changes when the contract provides discipline valued by the worker.
    \item The empirical design must separate commitment from coercion and sorting.
\end{itemize}
\end{frame}

\begin{frame}{Theory-to-applied translation}
\small
\begin{enumerate}
    \item Hidden object: effort, type, commitment, evaluator favoritism, multitasking, peer effects, or relational enforcement.
    \item Observable margin: output, quality, attendance, quits, task mix, ratings, promotions, wages, or contract take-up.
    \item Variation: contract change, field experiment, menu, probation rule, evaluator assignment, promotion threshold, or platform policy.
    \item Mechanism threats: incentives versus sorting, effort versus ability, quantity versus quality, output versus gaming.
    \item Welfare object: productivity, worker risk, retention, inequality, organizational output, client welfare, or worker surplus.
\end{enumerate}
\end{frame}

\begin{frame}{Frontier methods for hidden objects}
\small
\begin{tabular}{p{0.30\linewidth} p{0.56\linewidth}}
\toprule
Method & Best use \\
\midrule
Output admin data & high-frequency productivity, quality, pay, and tenure \\
Contract-change studies & real incentive, probation, bonus, or schedule redesigns \\
Randomized experiments & pay formula, framing, goals, or contract menus \\
Evaluator designs & favoritism, influence activities, subjective ratings \\
Surveys linked to admin data & commitment, beliefs, fairness, risk, self-control \\
Digital traces and attendance & effort timing, discipline, time use, platform activity \\
Structural estimation & latent effort, type, commitment, and counterfactual contracts \\
\bottomrule
\end{tabular}
\end{frame}

\begin{frame}{Research Lab design}
\begin{itemize}
    \item Primary anchor: performance pay and productivity as a contract-change labor design.
    \item Challenge anchor: subjective evaluation and influence activities; optional commitment path through self-control at work.
    \item Reproduce: show higher output after performance pay in a synthetic teaching panel.
    \item Diagnose: decompose output change into within-worker incentive response and workforce-composition change.
    \item Transfer: redesign the workflow for subjective evaluation or commitment contracts.
    \item Deliverable: hidden object, margin, variation, identification strategy, welfare object, and main threat.
\end{itemize}
\end{frame}

\begin{frame}{Takeaways}
\begin{itemize}
    \item Contracts are labor-market institutions: they allocate risk, elicit effort, reveal type, and govern retention.
    \item Principal-agent logic is most useful when paired with a real workplace signal and a real empirical design.
    \item Screening, effort incentives, multitasking, subjectivity, and relational enforcement are distinct mechanisms.
    \item The frontier contribution is measurement: make hidden effort, type, commitment, and evaluator behavior empirically legible.
\end{itemize}
\end{frame}

\end{document}
